\chapter{均衡（EQ）}
\label{chap:eq}

均衡不是“把声音调好听”的万能旋钮，而是\textbf{改变不同频率的相对电平}。同一段素材，削掉 300 Hz 的堆积和抬高 3 kHz 的亮度，听感会完全不同；但两者都不能代替正确的电平、声像和演奏。本章先建立频段地图，再讲夜半的三段塑形 EQ 能做什么、不能做什么。

\section{频段地图与掩蔽}

人耳可听范围约 20 Hz–20 kHz。混音里常把它分成几段看，但分段不是硬边界，相邻频段始终互相影响：

\begin{center}
\begin{tikzpicture}[font=\scriptsize, every node/.style={align=center}]
  \foreach \i/\lab in {0/{20–60\\次低频}, 1/{60–250\\低频}, 2/{250–500\\低中频}, 3/{0.5–2k\\中频}, 4/{2–5k\\presence}, 5/{5–10k\\高频}, 6/{10k+\\空气感}} {
    \fill[black!12] (\i*2.3,0) rectangle ++(2.3,1.2);
    \draw (\i*2.3,0) rectangle ++(2.3,1.2);
    \node at (\i*2.3+1.15,0.6) {\lab};
  }
  \draw[-{Latex[length=2mm]}] (0,-0.35) -- (16.1,-0.35) node[right] {频率升高};
\end{tikzpicture}
\end{center}

低频给重量与根基，低中频给 body，中频承载旋律与语言，中高区决定“ presence/清晰度”，高频决定空气感。两个音源在同一频段都很响时会发生\textbf{掩蔽}：人耳更注意较响的那个，另一个即使电平不低也会“消失”。EQ 的第一价值是减少这种争夺，不是给每轨加一个好看的曲线。

\section{先减法，再加法}

\begin{enumerate}
  \item \textbf{先找多余的}：扫频时先大幅提升一个窄段，找到“变浑/变刺”的位置，再反过来衰减。提升只是探针，最终曲线通常以衰减为主。
  \item \textbf{再补缺失的}：只有在减法之后仍缺重量或空气感时，才小幅提升；提升会同时抬高该段所有声音，包括噪声与齿音。
  \item \textbf{每步都回听整体}：单轨独奏好听不等于在整体里成立。每次改动后回到全混，检查其他元素是否被推远。
\end{enumerate}

扫频的常用做法：把中频增益临时推到 +9–12 dB，缓慢移动中心频率，听到“最难受”的点就停；然后把增益拉回 0 并向下衰减 2–4 dB。夜半当前没有独立 Q 旋钮，新装 EQ 的中频 Q 为默认 0.9；曲线拖动只改增益与中心频率，不改 Q。所以“窄扫”在本机架上做不到像专用均衡器那样收窄，探针会比理想状态宽一些，判断时要以整体听感为准，不要死盯一个频点。

\section{夜半的三段塑形 EQ}

在设备机架点“装 EQ”，会在当前轨的设备链上插入一台三段塑形 EQ。它由三级串联的 RBJ biquad 组成：低架、中频峰值、高架。某一段增益为 0 dB 时，该段在数学上是恒等变换；三段都为 0 dB 时，整台 EQ 是真正的旁通，不是“接近旁通”。

\begin{table}[htbp]
\centering
\small
\begin{tabular}{@{}llll@{}}
\hline
参数 & 作用 & 范围 & 默认 \\
\hline
\texttt{eq\_low\_gain} & 低架增益 & $-18$–$+18$ dB & 0 dB \\
\texttt{eq\_low\_freq} & 低架转折 & 20–1000 Hz & 200 Hz \\
\texttt{eq\_mid\_gain} & 中频峰值增益 & $-18$–$+18$ dB & 0 dB \\
\texttt{eq\_mid\_freq} & 中频中心 & 200–8000 Hz & 1000 Hz \\
\texttt{eq\_mid\_q} & 中频品质因数 & 0.3–6 & 0.9 \\
\texttt{eq\_high\_gain} & 高架增益 & $-18$–$+18$ dB & 0 dB \\
\texttt{eq\_high\_freq} & 高架转折 & 1000–16000 Hz & 4000 Hz \\
\hline
\end{tabular}
\caption{夜半三段塑形 EQ 的参数契约}
\label{tab:eq-params}
\end{table}

设备机架中央的频响图是\textbf{真实参数}算出的响应，不是装饰曲线。横轴为 20 Hz–20 kHz 对数刻度，纵轴为 $\pm18$ dB，中线是 0 dB。图上的无障碍标签仍写着“4 段 EQ 频响曲线”，但当前代码契约与实现都是三段：低架、一个中峰、高架；没有第四段，也没有可拖出的第四个点。

\screenshot{device-rack-eq-curve}{当前轨已装 EQ，设备机架中央显示频响曲线，低架增益滑块可见，曲线不是平线。}

\section{在曲线上拖动}

频响图不只是显示，它是一块可拖动区域。拖动位置会同时写两个参数：本段的增益和本段的频率。

\begin{table}[htbp]
\centering
\small
\begin{tabular}{@{}lll@{}}
\hline
横向位置 & 写入的增益 & 同时写入的频率 \\
\hline
左 1/3（$x<0.333$） & \texttt{eq\_low\_gain} & \texttt{eq\_low\_freq} \\
中 1/3（$0.333\le x<0.667$） & \texttt{eq\_mid\_gain} & \texttt{eq\_mid\_freq} \\
右 1/3（$x\ge0.667$） & \texttt{eq\_high\_gain} & \texttt{eq\_high\_freq} \\
\hline
\end{tabular}
\caption{频响图拖动落点}
\label{tab:eq-drag}
\end{table}

纵向以中线为 0 dB，向上向下各到 $\pm18$ dB；增益按 0.1 dB 取整。横向按 $20\times10^{3x}$ Hz 求频率，再夹到该段的合法范围。也就是说，曲线最左端不会低于 20 Hz，低段拖到中段边界时频率最多 1000 Hz；中段最多 8000 Hz；高段最多 16000 Hz。曲线拖动\textbf{不改} \texttt{eq\_mid\_q}，Q 保持默认或文件里已有的值。

低架增益还有一个独立的滑块（\texttt{eq\_low\_gain}），它只改增益、不改转折频率。日常修低频堆积时用滑块更直接；需要同时移动低架转折点时再用曲线左段。主总线也可以装 EQ，并提供同样的低架增益滑块；主总线的处理影响整个工程，动之前先把展示轨的 EQ 调好。

\section{常见频段处理}

\begin{table}[htbp]
\centering
\small
\begin{tabular}{@{}p{0.18\linewidth}p{0.36\linewidth}p{0.38\linewidth}@{}}
\hline
频段 & 常见问题 & 处理方向 \\
\hline
20–60 Hz & 次低频，太多会闷、吃余量 & 高通或低架衰减；先确认监听能听到 \\
60–250 Hz &  bass 与 kick 的根基，堆积会“轰” & 低架小幅衰减，或让各轨错开 \\
250–500 Hz & 低中频“纸箱感” & 窄衰减 2–3 dB，别把 body 削空 \\
500 Hz–2 kHz & 旋律与语言的核心，过多会“电话音” & 先查编曲重叠，再小幅衰减 \\
2–5 kHz &  presence，过多会刺、齿音重 & 中峰小幅衰减；提升要非常克制 \\
5–10 kHz & 清晰度与齿音 & 高架微调，配合齿音处理 \\
10 kHz+ & 空气感 & 高架小幅提升；噪声也会一起抬高 \\
\hline
\end{tabular}
\caption{常见频段问题与处理方向}
\label{tab:eq-bands}
\end{table}

这些数字是起点，不是处方。真正要听的是“这个音源在整体里是否还需要占这么多空间”。如果两轨在同一段互相掩蔽，优先用声像、电平、编曲错位解决；EQ 只是其中一件工具。

\section{边界与常见错误}

\begin{itemize}
  \item \textbf{不要用 EQ 修削波}：EQ 提升会改变峰值，削波要在增益结构和动态处理里解决。
  \item \textbf{不要每轨都抬高频}：叠加后整体会刺；高频提升要按整体预算分配。
  \item \textbf{不要只看曲线}：曲线平不代表好听，曲线弯也不代表错；以全混回听为准。
  \item \textbf{多台 EQ 的曲线只画第一台}：设备机架的频响图跟随链上第一台带 EQ 参数的设备；链上装第二台 EQ 时，第二台的参数仍在处理，但图上不一定显示。
  \item \textbf{主总线 EQ 影响全局}：主总线上的低架增益会作用于整个工程，不要在单轨问题时先动主总线。
  \item \textbf{界面没有 Q 旋钮}：中频 Q 在参数契约里存在，但当前机架没有独立 Q 控件；需要更窄/更宽的频带时，现有界面做不到，不要假装拖动曲线能改 Q。
\end{itemize}

\section{帮助主题映射}

\begin{table}[htbp]
\centering
\small
\begin{tabular}{@{}p{0.40\linewidth}p{0.48\linewidth}@{}}
\hline
小节 & 应用内帮助锚点 \\
\hline
频段地图与掩蔽 & \texttt{eq-frequency-map} \\
先减法，再加法 & \texttt{eq-subtract-first} \\
夜半的三段塑形 EQ & \texttt{eq-three-band} \\
在曲线上拖动 & \texttt{eq-curve-drag} \\
常见频段处理 & \texttt{eq-common-bands} \\
边界与常见错误 & \texttt{eq-boundaries} \\
\hline
\end{tabular}
\caption{第 3 章帮助主题映射}
\end{table}
